% ODD+D description of the v4 detection-boundary model and experiment.
% Claim markers ODDxx refer to rows in claim-evidence.md, section
% "ODD appendix and model verification". Code locators are relative to
% studies/ic-2026-2027/src/ic_study/; protocol locators are relative to
% studies/ic-2026-2027/research/questions/exploration/2026-09-24-parallel/.
% No sentence in this file reports a gate outcome, AUC value or boundary.

This appendix describes the simulation model and the detection experiment built on it, following the second update of the ODD protocol \cite{grimm2020odd}. It also addresses the decision-making elements of the ODD+D extension \cite{muller2013oddd}. % ODD41
For numerical reproduction, this appendix describes the implementation at \code{37513841}, with each submodel naming its Python function. The frozen documents retain their original wording and authority; this account does not amend them. Where document and code differ, reproduce the frozen implementation using the readings and deviations already recorded in the docking ambiguities and integration decisions. These cover inclusive inversion and strict crossing claims (Q8), sustained boundaries regardless of reversal source (Q12), and main-curve validity without an eligibility exemption (Q13). G2 substitution follows the later amendment and code, as explained in Submodel~14. Methods and Section~\ref{sec:verification} describe these implementation readings. % ODD41

% ODD01
\paragraph{Identification.}
\emph{Model:} v4 detection-boundary successor (amendment v4.3), a linked successor of the frozen v1 protocol. \emph{Code revision:} Git commit \code{37513841}, whose record \code{freeze-v4.json} has status \code{frozen} and code-closure SHA-256 \code{db9d6ead}\allowbreak\code{b5019d69}\allowbreak\code{1a8c7370}\allowbreak\code{3d511d03}\allowbreak\code{37dc4983}\allowbreak\code{4bb219d1}\allowbreak\code{68a27693}\allowbreak\code{df0e0acd}. \emph{Protocol:} \code{detection-boundary-amendment-v4.md}, the decision specification D1--D11, revision R1--R8 and two review-decision records, all hash-bound in that freeze record. \emph{Run manifest:} the report manifest accompanies the exported results. \emph{Validity domain:} synthetic eleven-wave panels with the declared geometry, laws, response channel and discriminator. No parameter or panel of the simulation is estimated from survey microdata; the separate SHP reliability anchor is described under Input data.

\subsection{Purpose and patterns}

% ODD02
\paragraph{Purpose.}
The model asks how much nominal partner drift makes synthetic opinion panels distinguishable from matched panels generated without interaction, when both pass through the same observation channel. In one arm, individuals update a latent opinion by a rank-transported autoregression whose innovation includes a standardized, bounded-confidence-shaped partner drift with repulsion on a fixed random network, with nominal strength $a$. In the other arm, a null family produces persistence and updating without any network. Both arms are observed through the same attrition process, response-error channel with signal share $r$ and eleven-category ordinal allocation, except that in comparisons against null C the null arm alone receives heteroscedastic error and heaping (Submodel~5). The primary estimand is the whole curve $a^*(r)$ over nine values of $r$: the smallest strength on a fixed grid at which a fixed linear discriminant reaches an area under the ROC curve (AUC) of 0.70 against the confirmatory null B. The model is a demonstration and measurement device for a synthetic detection limit. It is not calibrated to predict empirical opinion change, and it does not identify causal social influence: no exposure network is observed in the panels it imitates \cite{manski1993,shalizi2011}.

% ODD03
\paragraph{Patterns.}
The patterns are synthetic design targets, fixed before any v4 analysis, not observations to be reproduced. (P1) Panel geometry: 10,626 people in two-person households, 11 annual waves, entry of 7,084 people at wave 0 and 3,542 at wave 5, and household retention 0.89 per wave. This approximates, but is not reconstructed from, the Swiss Household Panel trust series used by Gestefeld and Lorenz \cite{gestefeld2023}. (P2) Matched short-lag persistence: each null targets the mean observed-scale lag-1 correlation (and lag-2 for two-parameter nulls) of the interaction arm in the same cell. Matching follows the root and constrained-face rules of Submodel~8 and independent validation of Submodel~9. (P3) Identical observed marginals in both arms under a common mask, which the wave-wise rank allocation guarantees by construction. (P4) Pathwise identity at zero strength: at $\phi=0.85$, zero mixture weight and $\omega=0$, nulls A, B, D and E coincide exactly with the interaction arm at $a=0$ on shared seeds. The model is not designed to reproduce the empirical category frequencies of the trust item or its empirical autocorrelation function, and neither is computed here. Whether the discriminator separates the arms is reported in Results.

\subsection{Entities, state variables and scales}

% ODD04
\paragraph{Individuals.}
There are $N=10{,}626$ individuals $i$ per panel. Each has a latent opinion $z_{it}$, a real number on the standard-normal rank scale that the transport $R_N$ (Submodel~1) produces. Each also has a household index that determines attrition. Arm-specific fixed attributes are:
\begin{itemize}
\item eight out-partners $\pi_1(i),\ldots,\pi_8(i)$ in the interaction arm;
\item a trait $T_i\sim N(0,1)$ in null B;
\item a stayer indicator $s_i\sim\mathrm{Bernoulli}(\pi)$ in null E;
\item a response-style component $u_i\sim N(0,1)$ in the persistent-style sensitivity.
\end{itemize}
Observed quantities per person-wave are a response indicator $m_{it}\in\{0,1\}$, a continuous response $w_{it}$ and an ordinal response $y_{it}\in\{0,\ldots,10\}$, defined only when $m_{it}=1$. Latent states are never observed by the discriminator, except in the declared latent control channels.

% ODD05
\paragraph{Environment and collectives.}
The interaction network is a fixed directed multigraph formed by eight independent uniform permutations of $\{1,\ldots,N\}$. Each person has eight out-edges and eight in-edges counting multiplicity; self-loops and repeated partners are allowed. It has no spatial embedding and is never rewired. Households of two people are the only collectives. They matter only for attrition, not for opinion updating.

% ODD06
\paragraph{Parameters and scales.}
Time is measured in waves $t=0,\ldots,10$, interpreted as years. The main design has one synchronous update step per wave; a sensitivity has 24 steps per wave, observed only at the end of the wave. The fixed parameters are:
\begin{itemize}
\item persistence $\rho=0.85$ per step;
\item repulsion $b=0.25$, with an attraction-only sensitivity $b=0$;
\item confidence distance 1 on the latent scale;
\item anchoring $\omega=0$, with sensitivities $\omega\in\{0.25,0.50,0.75,1\}$.
\end{itemize}
The experimental factors are the strength $a\in G=\{0,0.001,0.003,0.01,0.03,0.10,0.30,0.60,0.90\}$ and the signal share $r\in R=\{0.50,0.55,\ldots,0.90\}$ of the continuous response channel. The parameters $a$ and $r$ are dimensionless shares. Null parameters are fitted (Submodel~8), not set: $\phi,\rho_0\in[0,0.999]$ and $\eta\in\{\lambda,\pi,\gamma\}\in[0,0.8]$. These boxes define the declared contrast, not the full mathematical persistence families (Supplementary Table~\ref{tab:design-conventions}).

\subsection{Process overview and scheduling}

% ODD07
One panel realization proceeds as follows (\code{simulation/detection_v4.py:PanelSource._profiles}):
\begin{enumerate}
\item Initialize latent opinions $z_{\cdot0}=R_N(\varepsilon_0)$. In the interaction arm with $a>0$, draw the network.
\item For each wave $t=1,\ldots,10$ and each of $S\in\{1,24\}$ steps:
  \begin{enumerate}
  \item compute every individual's bounded drift from the current states of its partners;
  \item compute all updates from the same pre-step state vector (synchronous update);
  \item apply the rank transport $R_N$ to the whole population, which commits the new state.
  \end{enumerate}
  Null arms replace (a)--(b) by their own recursion (Submodel~4) and apply $R_N$ at every step.
\item Draw the household attrition mask for all waves. It is independent of opinions, and latent states evolve for everyone, observed or not.
\item Generate the continuous response $w_{it}$ and allocate ordinal categories wave by wave among respondents.
\item Compute the panel's summary features (Submodel~6).
\end{enumerate}
Opinion updating stops after wave 10. There is no convergence criterion. Tied values are ordered by index (stable sort) in $R_N$. The detection experiment (Submodels 7--16) repeats this process over thousands of seeded panels per cell.

\subsection{Design concepts}

% ODD08
\paragraph{Basic principles.}
The interaction arm combines similarity-biased assimilation with repulsion beyond a distance, one of the influence classes distinguished by Flache et al.\ \cite{flache2017}, on a fixed random network. It is written as a persistent autoregression whose innovation is split between standardized partner drift and fresh noise. Because $R_N$ restores the same marginal at every step, the model keeps only the local shape of the bounded-confidence kernel, not clustering or polarization in the population marginal. The null families encode persistence and updating without interaction:
\begin{itemize}
\item rank-transported autoregression (A);
\item durable trait plus autoregression (B);
\item state-dependent persistence (D);
\item a mover--stayer mixture (E);
\item autoregression with a different error channel (C).
\end{itemize}
Measurement follows a classical true-score-plus-error channel before discretization. Detection is framed as two-sample classification of whole panels by a fixed linear discriminant. For each successful validated fit and training realization, the Bayes comparator knows the same $r$ and fitted null law. The trained discriminator's population AUC cannot exceed that conditional optimum; taking expectations under the same successful-fit procedure preserves the inequality. This is an average of conditional optima, not the optimum against a marginal mixture of fitted nulls. Non-attainment therefore does not establish indistinguishability by any procedure. The bound concerns population quantities, not estimated AUCs or inverted intervals, and it is probably loose, because an optimal procedure could use every person's trajectory rather than three pooled correlations.

% ODD09
\paragraph{Emergence.}
The lag profile of within-person correlations beyond lag 2, and in particular their decay at lags 3--5, arises from the interaction of drift, persistence, rank transport and the observation channel. The discriminator uses these lags. Marginal distributions are imposed by transport, and short-lag persistence is imposed on the nulls by matching. Neither is emergent.

% ODD10
\paragraph{Adaptation, objectives, learning and prediction.}
None apply. Individuals follow fixed update rules with no goals, memory beyond the current state (and the initial state under anchoring or, for stayers, in null E), learning or forecasting. In ODD+D terms there is no individual decision-making beyond this stochastic update rule \cite{muller2013oddd}. The only fitted components belong to the observer: null parameters and the discriminant. They are not agent behaviour.

% ODD11
\paragraph{Sensing and interaction.}
At each step an individual reads the current, error-free latent states of its eight partners. Interaction is direct and one-directional along the fixed edges. It enters only through the partner-averaged drift. The distance threshold compares latent differences on the unit-variance rank scale.

% ODD12
\paragraph{Heterogeneity and stochasticity.}
Rules are homogeneous. Individuals differ in initial state, network position and, depending on the arm, in $T_i$, $s_i$ or $u_i$. Randomness enters through:
\begin{itemize}
\item initial states and per-step innovations;
\item network permutations;
\item attrition;
\item response error, persistent style and heaping;
\item null traits and stayer indicators.
\end{itemize}
Every stream is a NumPy \code{default_rng} generator seeded from a SHA-256 seed material (Submodel~15). This allows exact replay and common random numbers where they are declared.

% ODD13
\paragraph{Collectives.}
Two-person households share entry and attrition. No group-level opinion, identity or decision exists.

% ODD14
\paragraph{Observation.}
The unit of observation and classification is one complete panel, not a person or a person-wave pair. Latent states and networks are regenerated from seeds and are not stored; per-panel summary features and each cell's fitted null parameters are kept in the sealed store, and the parameters are released only with complete units under the linked release rule in Submodel~16. % ODD14 RC43
The analysis uses the ordinal responses of respondents, summarized as pooled lag correlations (Submodel~6), and, for declared controls, masked latent values. Each experimental cell records:
\begin{itemize}
\item solver statuses and evaluation counts;
\item validation residuals;
\item test scores, nuisance draws and tie rates.
\end{itemize}

\subsection{Initialization}

% ODD15
All panels start from $x_0=R_N(\varepsilon_0)$ with $\varepsilon_0\sim N(0,I_N)$, without burn-in; this is $z_{\cdot0}$ in the interaction arm and in nulls A, D and E, whereas null B starts from $R_N(\sqrt\lambda\,T+\sqrt{1-\lambda}\,x_0)$. Adjustment of the states after the network is drawn is therefore part of the partner-drift signal. The network is drawn once per panel. The interaction arm uses no network at $a=0$. Arms differ only in $a$, the null family and its fitted parameters, and the declared channel or sensitivity. Initial states are synthetic, not data-based. Seed roles and indices (Submodel~15) determine which arms share initial draws.

\subsection{Input data}

% ODD16
The model uses no external time-varying input and no survey microdata. The Swiss Household Panel appears in two roles only. First, its trust-in-federal-government item \code{P$$P04} (1999--2009) motivates the synthetic geometry: Gestefeld and Lorenz report 10,627 people, and $N$ is rounded to 10,626 for divisibility by six \cite{gestefeld2023}. Second, it is the object of a separately frozen, secondary preregistration that estimates an empirical reliability anchor. That anchor was computed locally, on the first author's computer, from SHP microdata (version 12.0 \cite{shpgroup2026data}) obtained under the SHP data-use contract, which does not allow redistribution; only aggregate estimates, standard errors, fit statistics and counts left the restricted environment, and no microdata enter the simulation or the repository.

% ODD17
\paragraph{Empirical reliability anchor (secondary, preregistered).}
The anchor preregistration was frozen before any SHP microdata were accessed. It specifies a Gaussian FIML quasi-simplex model $Y_{it}=\mu_t+Z_{it}+\epsilon_{it}$ with $Z_t=\phi Z_{t-1}+\eta_t$ over the eleven waves, under MAR. The anchor is the 2004 reliability $r_{2004}=\mathrm{Var}(Z_t)/(\mathrm{Var}(Z_t)+\mathrm{Var}(\epsilon_t))$, with a Wald 95\% interval. Heise and Wiley--Wiley three-wave estimates serve as specification sensitivities \cite{heise1969,wiley1970}. When available, the anchor is displayed as a point over the frozen grid $R$. It changes no grid, target, null, band, censoring or gate rule, and $a^*(r)$ results may not be used to choose its specification. The anchor value, or its unavailability, is reported in Results. The anchor treats the 0--10 response as continuous. The simulation's $r$ is instead the signal share of a continuous channel before rank-based discretization. The two quantities are therefore comparable only approximately.
% ODD39
Successor 1 was frozen on 28 September 2026 at 07:51:17 UTC, with code revision \code{17a0dbed} bound in \code{shp-reliability-successor-1-freeze.json}, before SHP microdata access. % ODD39 RC44
It adds a trait--state--error sensitivity \cite{kenny1995tse}, reported alongside the unchanged quasi-simplex anchor and never selected by fit; a descriptive crosswalk of the anchor estimators on synthetic panels of known signal share; and a description of the 2004 refreshment wave. % ODD39 RC45
The raw anchor, admissible STARTS sensitivity and synthetic compatibility set remain separate: the set depends on the synthetic law and uses minimum--maximum ranges over ten replicates, not confidence intervals or an uncertainty-propagating calibration of $r$. % ODD39 RC46

\subsection{Submodels}

Submodels 1--6 define the generative model. Submodels 7--16 define the detection experiment performed with it. Following Grimm et al.\ \cite{grimm2020odd}, the latter are simulation-experiment design rather than model description, but they are specified here with the same precision. Constants are those of the frozen code.

\subsubsection*{1. Normal-rank transport $R_N$}
% ODD18
For $v\in\mathbb R^N$, let $k_i\in\{0,\ldots,N-1\}$ be the rank of $v_i$ under a stable sort, so ties are ordered by index. Define the grid $\tilde q_k=\Phi^{-1}((k+\tfrac12)/N)$ and its standardization $q_k=\tilde q_k/\mathrm{sd}(\tilde q)$, using the population standard deviation. Then
\begin{equation*}
R_N(v)_i=q_{k_i}.
\end{equation*}
$R_N$ preserves ranks and returns the same mean-zero, unit-variance marginal at every step. Code: \code{models/anchored_panel.py:normal_ranks} and the byte-identical fast path \code{simulation/amendment_a.py:fast_ranks}.

\subsubsection*{2. Panel geometry and attrition}
% ODD19
Persons $2h$ and $2h+1$ form household $h=0,\ldots,N/2-1$. Households $h<N/3$ (7,084 people) enter at wave 0; the rest (3,542 people) enter at wave 5. With $A_{h,-1}=0$ and $U_{ht}\sim U(0,1)$,
\begin{align*}
A_{ht}&=\bigl(A_{h,t-1}\wedge U_{ht}<0.89\bigr)\vee\mathbf 1[e_h=t],\\
m_{it}&=A_{h(i)t},
\end{align*}
where $e_h$ is the entry wave. Exit is permanent and independent of opinions. Code: \code{models/anchored_panel.py:panel_mask}, called with retention 0.89 from \code{PanelSource._profiles}.

\subsubsection*{3. Influence process: bounded confidence with repulsion}
% ODD20
Drift is computed from the current state $z$:
\begin{align*}
D_i&=\frac18\sum_{k=1}^8 f\bigl(z_{\pi_k(i)}-z_i\bigr),\\
f(g)&=\begin{cases}g,&|g|\le1\\-b\,g,&|g|>1\end{cases}
\end{align*}
It is then standardized across the population: $d_i=(D_i-\bar D)/\mathrm{sd}(D)$. A zero drift standard deviation at $a>0$ raises an error instead of substituting another kernel. The update is
\begin{align*}
u_i&=\rho z_i+\sqrt{1-\rho^2}\bigl(\sqrt a\,d_i+\sqrt{1-a}\,\varepsilon_i\bigr),\\
z'&=R_N\bigl((1-\omega)u+\omega z_{\cdot0}\bigr),
\end{align*}
with $\rho=0.85$, $b=0.25$, $\omega=0$ and $\varepsilon_i\sim N(0,1)$. The strength $a\in[0,1)$ is the nominal share of the pre-transport innovation variance given to drift. It is not the realized share of latent change. At $a=0$ the recursion is $z'=R_N(\rho z+\sqrt{1-\rho^2}\varepsilon)$. The sensitivities are $b=0$, $\omega\in\{0.25,0.5,0.75,1\}$ (with $\omega=1$ an explicit degeneracy in which every wave repeats $z_{\cdot0}$ whatever $a$ is), and $S=24$ steps per wave with the same per-step coefficients. Code: \code{models/detection_panel_v4.py:influence}, which calls \code{models/detection_panel.py:influence_update} and \code{bounded_drift}.

\subsubsection*{4. Null families}
% ODD21
All nulls share the finite-population transport of the interaction arm. Let $X_\phi$ be the zero-strength recursion with the same number of transports per wave as the alternative:
\begin{align*}
x_0&=R_N(\varepsilon_0),\\
x_t&=R_N\bigl(\phi x_{t-1}+\sqrt{1-\phi^2}\,\varepsilon_t\bigr).
\end{align*}
The families are (code: \code{models/detection_panel_v4.py:null}):
\begin{itemize}
\item \textbf{B (confirmatory; trait plus AR).} $z_t=R_N(\sqrt\lambda\,T+\sqrt{1-\lambda}\,x_t)$, with $T_i\sim N(0,1)$ i.i.d., independent of $x$ and not rank-transported.
\item \textbf{A.} B at $\lambda=0$, i.e.\ $z=X_\phi$.
\item \textbf{E (mover--stayer).} $z_t=R_N(s_i x_0+(1-s_i)x_t)$ with $s_i\sim\mathrm{Bernoulli}(\pi)$.
\item \textbf{D (state-dependent persistence).} $z_0=R_N(\varepsilon_0)$ and $z'=R_N(\rho(z)z+\sqrt{1-\rho(z)^2}\,\varepsilon)$ at each step, with $\rho(z)=\rho_0+\gamma(1-\rho_0)\bigl(1-e^{-z^2/2}\bigr)$. At $\gamma=0$, D equals $X_{\rho_0}$.
\item \textbf{C.} The latent law of A, observed through the heteroscedastic-error and heaping channel of Submodel~5.
\end{itemize}
The weights $\lambda$ and $\pi$ are nominal pre-transport weights. The Gaussian covariance identities $\lambda+(1-\lambda)\phi^k$ and $\pi+(1-\pi)\phi^k$ of v1 are no longer claimed for these transported laws. At $\phi=0.85$ and $\lambda=\pi=\gamma=0$, nulls A, B, D and E coincide pathwise with the interaction arm at $a=0$ when they share state seeds. No identity is claimed for C or for $\omega>0$.

\subsubsection*{5. Response channel, reliability and ordinal allocation}
% ODD22
The continuous response is
\begin{equation*}
w_{it}=\sqrt{r_{it}}\,z_{it}+\sqrt{1-r_{it}}\;\sigma_{it}\,e_{it},
\end{equation*}
with $e_{it}=\sqrt{1-c}\,v_{it}+\sqrt c\,u_i$, where $v_{it},u_i\sim N(0,1)$ are independent of states, network and mask. In the main channel $c=0$, $\sigma\equiv1$ and $r_{it}=r$. Response error means transient variation that is not true attitude; true shocks belong to $z$. The parameter $r$ is thus the signal share of the standardized continuous channel before discretization, a sensitivity parameter rather than an estimated reliability. The sensitivities are:
\begin{itemize}
\item persistent style, $c\in\{0.25,0.50\}$;
\item panel conditioning, $r_{it}=\min(0.95,\,r+\delta\,\tau_{it}/10)$ with $\delta\in\{0.10,0.20\}$, where $\tau_{it}$ counts the person's earlier completed interviews and is 0 at first interview;
\item the exploratory combination $c=0.50,\ \delta=0.20$.
\end{itemize}
For the null arm of C only, error is heteroscedastic. Let $\tilde\sigma(z)=1+2e^{-z^2/2}$. Then $\sigma_{it}=\tilde\sigma(z_{it})\big/\bigl(\overline{\tilde\sigma^2}\bigr)^{1/2}$, where the mean runs over all $11\times N$ person-waves of the panel.

% ODD23
Ordinal allocation is carried out wave by wave among the $n_t$ respondents. Let $p_\ell=(1-\mu)/11+\mu\,\mathbf 1[\ell=5]$ for $\ell=0,\ldots,10$, and $C_j=\sum_{\ell\le j}p_\ell$. The respondent of rank $k\in\{0,\ldots,n_t-1\}$ in $w_{\cdot t}$ receives
\begin{equation*}
y=\bigl|\{j\in\{0,\ldots,9\}:C_j<(k+\tfrac12)/n_t\}\bigr|.
\end{equation*}
In the main design $\mu=0$, giving eleven approximately equiprobable categories. In comparisons against C, both arms use $\mu=0.1$, a common target marginal. In the C null arm only, each respondent's $w_{it}$ is set to 0 with probability 0.10 before ranking (heaping), and ties are broken by an independent random permutation. The wave-wise rank allocation gives both arms identical observed marginals under a common mask. Because it uses the whole wave, it removes marginal information. Code: \code{models/detection_panel.py:measurement}, \code{ordinal_channel} and \code{tenure_reliability}.
The effective thresholds depend on the respondents' distribution in each wave, inducing dependence among their categories. A person's category can therefore change with respondent composition even if the continuous response stays fixed. The same allocation operates in both arms, with the panel still the classification unit. The boundary refers to this relative discretization and does not automatically apply to a 0--10 scale with fixed individual thresholds. % ODD23

\subsubsection*{6. Features and channels}
% ODD24
For a panel of values $x_{it}$, with missing entries excluded, define
\begin{equation*}
L_k=\mathrm{corr}\bigl(\{(x_{it},x_{i,t+k})\}\bigr),\quad k=1,\ldots,5.
\end{equation*}
The correlation is the Pearson correlation over all pairs of the panel in which both waves are observed, pooled over persons and waves. Matching uses $(L_1,L_2)$; the discriminator uses $(L_3,L_4,L_5)$. The channels are:
\begin{itemize}
\item \emph{observed}: $y$;
\item \emph{latent}: masked $z$;
\item \emph{ordinal\_one}: allocation of masked $z$ without error or heaping, i.e.\ $r=1$ with the family's common marginal;
\item \emph{latent\_rank}: per-wave mid-ranks of masked $z$, scaled to $(0,1)$.
\end{itemize}
The mobility summaries apply to the observed, latent and ordinal\_one channels. They replace $(L_3,L_4,L_5)$ by the mean absolute one-wave change $|x_{i,t+1}-x_{it}|$ within each tercile of the wave-$t$ mid-rank position, with terciles $[0,\tfrac13)$, $[\tfrac13,\tfrac23)$ and $[\tfrac23,1]$. Code: \code{simulation/detection_v4.py:selected_features}, \code{analysis/detection_boundary.py:lag_profile} and \code{midranks}.

\subsubsection*{7. Pre-discrimination feasibility screen}
% ODD25
The screen applies before matching, for observed channels only. With continuous i.i.d.\ error, $|\mathrm{corr}(w_t,w_s)|\le\sqrt{r_tr_s}$. With persistent style, the bound becomes $\sqrt{r_tr_s}+c\sqrt{(1-r_t)(1-r_s)}$. The code evaluates this bound at $r_t=r_s=\bar r$, where $\bar r=r$, or $\bar r=\min(0.95,r+\delta)$ under conditioning. A cell whose mean target $L_1$ exceeds the bound in absolute value is not matched. It becomes a gap, never counted as non-attainment. The continuous bound is a conservative operational rule for ordinal data, not an ordinal impossibility theorem. Code: \code{analysis/detection_boundary.py:feasible_lag1}, applied in \code{simulation/detection_v4.py:fit_cell}.

\subsubsection*{8. Transport-matched fitting: coupling, profile solver and budget}
% ODD26
All null arms receive the same finite-population transport as the interaction arm (Submodel~4). The null parameters $\theta=(\phi,\eta)$ are fitted to the interaction arm's observed $(L_1,L_2)$. In the matching phase, realization $j=0,\ldots,M-1$ of both arms uses index $2j+1$, so target and null share the state, mask and error streams. Components without a counterpart ($T$, $s$) use a separate \code{null-state} role, and the interaction arm's network uses its own stream. Let $g_j(\theta)$ be the null-minus-target difference in $(L_1,L_2)$ for pair $j$. Then
\begin{align*}
\bar g_k(\theta)&=\frac1M\sum_j g_{jk}(\theta),\\
\mathrm{SE}_{\mathrm{fit},k}&=\frac{\mathrm{sd}_j(g_{jk})}{\sqrt M},
\end{align*}
with $M=1{,}024$ on the first attempt.
Here $\mathrm{Var}(g_{jk})=\mathrm{Var}(L_{k,\mathrm{null}})+\mathrm{Var}(L_{k,\mathrm{target}})-2\mathrm{Cov}(L_{k,\mathrm{null}},L_{k,\mathrm{target}})$. The factor 32 stated in the Supplementary Material (\hyperref[supp:matching-scope]{scope of matching precision and residual tolerance}) compares an individual target statistic with its iid mean, not with this coupled residual standard error. % ODD26

% ODD27
The solver is a profile solver. For fixed $\eta$ it solves $\bar g_1=0$ in $\phi$, and then it solves $\bar g_2=0$ in $\eta$ along that profile. Lag 1 has priority by definition. The search proceeds as follows:
\begin{itemize}
\item The anchor $(0.85,\eta_0)$ and forward probes of 0.01 in each coordinate give a finite-difference Jacobian. A linear prediction from it seeds the search.
\item Inner solves use bracketed secant steps with bisection fallback. They use a looser provisional tolerance, $\max(\tau_1,\min(5\times10^{-5},0.02|\bar g_2|))$, away from the outer root, and are tightened before acceptance.
\item The outer step is Illinois-modified regula falsi, with a bisection step every fourth iteration.
\item The $L_2$ residual must not decrease in $\eta$, beyond the larger of the two $L_2$ tolerances, between any two evaluated profile points; an apparent reversal is re-checked after refining the inner roots to tolerance $\min(10^{-6},0.1\tau_1)$ and, if it persists, gives \code{non-monotone-profile}.
\end{itemize}
The one-parameter families (A, C) and the control channels solve only lag 1. Controls hold $\eta$ at the value fitted to the corresponding observed channel. An interior solution is accepted if
\begin{equation*}
|\bar g_k|\le\tau_k=\min\!\bigl(5\times10^{-5},\max(10^{-6},0.25\,\mathrm{SE}_{\mathrm{fit},k})\bigr)
\end{equation*}
holds for every fitted moment after at least one real solver step. A converged point that the solver never stepped to is labelled \code{boundary}, never \code{interior}. The lower face $\eta=0$ is accepted as \code{boundary} when the $L_1$ root there has an $L_2$ residual $\ge-\tau_2$ and the profile re-evaluated at $\eta=0.01$ increases.
The tolerance bounds the fitting sample's mean residual for equations that meet it, not the population mismatch. The one-sided face rule can leave $L_2$ above root tolerance for profiled B, D, and E. These fits still require independent validation, and the centered nuisance variance does not automatically include that mismatch. % ODD27

% ODD28
The budget is 96 null evaluations actually executed per attempt, each of $M$ panels. Jacobian probes count toward it; cached repeats do not. The accepted statuses are:
\begin{itemize}
\item \code{anchor-exact}: at $a=0$ for A, B, D and E, every paired difference is exactly zero, with no real step required;
\item \code{interior};
\item \code{boundary}.
\end{itemize}
The remaining statuses (\code{no-L1-root}, \code{non-monotone-profile}, \code{ill-conditioned}, \code{numerical-limit}, \code{family-limit-lower/upper}, \code{budget}, and \code{dependency-gap} for controls whose observed fit failed) are gaps. They are estimator states, not scientific findings. Code: \code{analysis/detection_boundary_v4.py:profile_match}; orchestration in \code{simulation/detection_v4.py:fit_cell} and \code{pipeline}.

\subsubsection*{9. Independent matching validation and retry}
% ODD29
Every accepted fit is re-evaluated at $\hat\theta$ on 2,048 fresh coupled pairs, drawn in a separate validation namespace. It fails if, for any fitted moment $k$,
\begin{equation*}
|\mathrm{res}_k|>4\sqrt{\mathrm{SE}_{\mathrm{fit},k}^2+\mathrm{SE}_{\mathrm{val},k}^2}.
\end{equation*}
This includes $L_2$ for B, D and E fits at the face. After a validation failure, and only then, one retry refits with $M=4{,}096$ under a new 96-evaluation budget and a fresh 2,048-pair validation. A second failure is a gap, and both attempts are kept. Two quantities are reported as diagnostics, not acceptance rules: the equivalence statistic $|\mathrm{res}_k|+2.2414\,\mathrm{SE}_{\mathrm{val},k}$ compared with 0.0003, and a Mahalanobis amplification of the residual into $(L_3,L_4,L_5)$. Code: \code{analysis/detection_boundary_v4.py:validate_fit} and \code{amplification}; retry logic in \code{simulation/detection_v4.py:fit_cell}.

\subsubsection*{10. Discriminator and train/test/reference split}
% ODD30
The discriminator is linear discriminant analysis on $(L_3,L_4,L_5)$. It is trained on 1,000 networks $\times$ 2 events per class (2,000 panels per class). Features are standardized with the pooled training mean and standard deviation. The pooled within-class covariance $S$ (divisor $n-2$) is shrunk to $\Sigma=0.8S+0.2I$. The direction is $w=\Sigma^{-1}(\bar x_1-\bar x_0)$, and the score is
\begin{equation*}
s(x)=\Bigl(\tfrac{x-c}{\sigma}-\tfrac{\bar x_0+\bar x_1}2\Bigr)^{\!\top} w.
\end{equation*}
Equal priors apply, the orientation is fixed by the training class means, and larger scores indicate the interaction arm. There is no threshold or hyperparameter search. Test and reference sets each have 2,000 networks $\times$ 2 events per class. The train, test (\code{check}) and reference phases are disjoint, and the two classes are independent: the null uses index $2j$ and the interaction arm $2j+1$, and the two events of a network share network index $\lfloor\text{index}/4\rfloor$. Training always uses first-attempt training seeds, so strengths stay aligned when a cell used the retry. Code: \code{analysis/redesign.py:fit_discriminant}; phases in \code{simulation/detection_v4.py:pipeline}.

\subsubsection*{11. AUC criterion and boundaries}
% ODD31
The test AUC is the Mann--Whitney probability over all score pairs, with half credit for ties (\code{analysis/detection_boundary.py:auc}). For each validated fit and trained score, population AUC is the ranking probability for two independent new panels, one from each arm, independent of fitting, validation and training data. The target $\mathcal A(a,r)$ averages this probability over the declared procedure, including the permitted retry, conditional on successful validated fits. With target 0.70, the boundaries are
\begin{align*}
a^*(r)&=\min\{a\in G:\mathcal A(a,r)\ge0.70\},\\
a^*_s(r)&=\min\{a\in G:\min_{u\ge a}\mathcal A(u,r)\ge0.70\}.
\end{align*}
Targets 0.60 and 0.80 are secondary. There is no interpolation or monotone smoothing. Code: \code{analysis/detection_boundary.py:boundary} and \code{sustained_boundary}.

\subsubsection*{12. Uncertainty: score bootstrap and local nuisance propagation}
% ODD32
\emph{Score component.} Draw 2,000 cluster-bootstrap replicates of the fixed test scores, resampling networks separately in each class and keeping both events together. The same indices are used at every strength of a curve. $\mathrm{SE}_{\mathrm{scores}}$ is the bootstrap standard deviation.

\emph{Nuisance component.} At a non-anchor fit, draw $b=1,\ldots,2{,}000$ replicates:
\begin{enumerate}
\item Resample the $M$ matching differences with indices taken from 4,096 common uniforms, and take their uncentered mean $\bar g^{(b)}$.
\item Solve the linearized shift with the Jacobian $J$ of $(L_1,L_2)$ at $\hat\theta$. With $\eta$ free,
\begin{equation*}
\Delta\eta=-\frac{\bar g_2-J_{21}\bar g_1/J_{11}}{J_{22}-J_{21}J_{12}/J_{11}}.
\end{equation*}
For controls, $\Delta\eta$ is the observed channel's draw; for A and C it is 0. Clip $\eta$ to $[0,0.8]$, re-solve lag 1 with $\Delta\phi=-(\bar g_1+J_{12}\Delta\eta)/J_{11}$, and clip $\phi$ to $[0,0.999]$.
\item Shift the null training and reference features by $D\,\Delta\theta_b$, where $D$ is the common-random-number finite-difference Jacobian of the mean $(L_3,L_4,L_5)$ over the matching realizations. The step is $h=0.01$, taken inward at the upper face.
\item Resample training networks within each class, refit the discriminant and evaluate $\mathrm{AUC}_b$ on the fixed reference.
\end{enumerate}
$V_{\mathrm{nuis}}$ is the sample variance of $\mathrm{AUC}_b$, not divided by the number of draws. At \code{anchor-exact} fits, $\Delta\theta\equiv0$ and $V_{\mathrm{nuis}}=0$ analytically. The two components combine as $\mathrm{SE}^2=\mathrm{SE}_{\mathrm{scores}}^2+V_{\mathrm{nuis}}$. The mean shift, a fourth-moment MCSE and the projected fraction are reported. Code: \code{analysis/detection_boundary_v4.py:score_uncertainty}, \code{nuisance_draws}, \code{project_shift}, \code{mean_feature_jacobian} and \code{finite_differences}.

\subsubsection*{13. Band rule, null gate and censoring}
% ODD33
The band is computed per $r$ over the available strengths. Let $\epsilon=1/(2n)$ with $n=2{,}000$ networks per class. Clip $\hat p$ to $[\epsilon,1-\epsilon]$, set $\ell=\mathrm{logit}\,\hat p$ and $\mathrm{SE}_\ell=\max(\mathrm{SE},\epsilon)/(\hat p(1-\hat p))$. Replicates are
\begin{equation*}
p^*_b=\mathrm{clip}_{[0,1]}\bigl(\hat p+(A_b-\bar A)+\delta^{\mathrm{nuis}}_b\bigr),
\end{equation*}
where the same centered nuisance deviation is added at aligned strengths. The critical value is $c=\max(q_{0.95}(\max_a|t_b(a)|),\,2.7729)$, where $t_b(a)$ is the logit deviation of $p^*_b(a)$ from $\ell(a)$, studentized at the replicate. The band is $\mathrm{logit}^{-1}(\ell\pm c\,\mathrm{SE}_\ell)$, keeping the endpoint 0 or 1 when $\hat p$ is 0 or 1. It is simultaneous over available strengths within one comparison/$r$ unit, with no joint coverage declared across $r$, families or channels.

The null gate requires the $a=0$ cell to be available and its band to contain 0.5. Otherwise the unit is \emph{complete invalid} and its boundaries are suppressed. Intervals for $a^*$ and $a^*_s$ come from inclusive inversion: the first crossing of the upper band gives the lower end, and that of the lower band the upper end. A strict one-sided crossing claim requires the lower band to exceed 0.70. A gap blocks any boundary across it. The first boundary, and each end of its interval, are kept only below the first gap. The sustained boundary is kept only above the last gap. Otherwise the result is censored as ``gap at $a=\ldots$''. A unit of B that is not complete valid makes the main curve invalid without deleting the other units' boundaries. Code: \code{analysis/detection_boundary_v4.py:band} and \code{curve_report}, and \code{analysis/detection_boundary.py:boundary_report}; unit status in \code{simulation/v4_curve.py:complete_curve} and \code{_close_inventory}.
For C, the zero-strength rule determines operational eligibility without guaranteed arm identity, since only its null arm receives the altered channel. A failure can reflect channel asymmetry and, by itself, establishes neither an implementation defect nor influence; an eligible C boundary concerns latent law and channel together. % ODD33

\subsubsection*{14. Validation gates, failure allowances and stop rule}
% ODD34
All gates must pass, in the order G0, G1b, G1c, G2, G2c, before any confirmatory seed is used. G1a is the unit-test battery. Failure requires a documented successor, not a rerun. Code: \code{simulation/v4_plan.py:gate_tasks}, \code{simulation/v4_gate.py:run_task} and \code{aggregate}, and \code{simulation/v4_criteria.py}.
\begin{itemize}
\item \textbf{G0 (blind matching mask).} Scope: the 81 cells for B, A, C, D and E on the observed channel plus the three B controls (eight grids), with calibration seeds and common random numbers across $a$ and $r$. Only statuses, counts, residuals in SE units and validation outcomes are published. It passes if every zero-strength fit of A, B, D, E and the B controls is \code{anchor-exact}, no \code{budget} status occurs, and \code{numerical-limit} plus \code{ill-conditioned} statuses are at most 1\% of fits.
\item \textbf{G1a (identities).} The full-$N$ identity and zero-nuisance unit tests (Section ``Model verification'').
\item \textbf{G1b (band mechanics).} Five synthetic-score scenarios: Gaussian ICC 0, Gaussian ICC 0.5, $t_{5.2}$ clusters, rounded ties, and cross-slot dependence. Each has 2,000 networks $\times$ 2 events per class, 9 slots, equal classes and zero nuisance. There are 1,000 repetitions per scenario, and each scenario needs at least 935 joint coverages of 0.5.
\item \textbf{G1c (end-to-end at $a=0$).} 50 full pipelines at each $r\in\{0.50,0.70,0.90\}$. All 150 fits must be \code{anchor-exact}, with at most 3 pointwise failures at 2.7729 and $|\overline{\mathrm{AUC}}-0.5|\le3$ SE, where SE is the standard deviation of the 150 AUCs divided by $\sqrt{150}$. The tie rate is reported.
\item \textbf{G2 (procedure dispersion).} Cells $(a,r)=(0.003,0.70)$, $(0.01,0.70)$, $(0.03,0.70)$, $(0.10,0.70)$, $(0.30,0.70)$, $(0.03,0.50)$ and $(0.03,0.90)$. Each has 100 independent pipelines forming 50 disjoint pairs. The statistic is the variance of within-pair logit differences divided by the mean summed logit variance, per cell and pooled. It uses pairs without a gap, and a percentile bootstrap with 399 pair resamples gives 95\% intervals. Pass requires every cell's upper limit $\ge0.4976013738$ and lower limit $\le1.2879137798$, and the pooled interval to satisfy $0.8183809739$ and $1.0545065407$ in the same way. Each cell's median of $V_{\mathrm{nuis}}(h{=}0.005)/V_{\mathrm{nuis}}(h{=}0.01)$ must lie in $[0.8,1.25]$. Covariance-Jacobian corrections exceeding $0.1\sqrt{V_{\mathrm{nuis}}}$ are flagged without changing the criteria.
\item \textbf{G2c (local propagation versus full refits).} 100 refits at $(0.10,0.70)$ and 40 at $(0.03,0.70)$, each with new matching, validation and training. The null reference is regenerated at $\theta_{\mathrm{refit}}$ with the baseline's seeds. Two D5 baselines provide the comparisons. A comparison fails if the one-sided 95\% bootstrap upper bound of $V_{\mathrm{D5}}/\mathrm{Var}_{\mathrm{refit}}$ is below 0.8 or its lower bound exceeds 2.0.
\end{itemize}
% ODD35
The end-to-end G2/G2c tasks use B on the observed channel, not every family or sensitivity. The \hyperref[supp:gate-calibration]{supplementary calibration description} gives the artificial model used to calibrate their criteria and its limits. % ODD34
Technical failures (gap statuses) remain in the planned denominators. The allowances are:
\begin{itemize}
\item G2: at most 5 of 100 per cell and $\lceil0.02\times\text{pipelines}\rceil$ pooled, which is 14 of 700;
\item G2c: at most 3 of 100 and 2 of 40;
\item G1c: none.
\end{itemize}
The thresholds come from artificial calibration with a declared 1\% failure rate (\code{criteria-v43.json}). Its targets are pass probability at least 0.95 with correct standard errors and at most 0.20 with standard errors understated by 20\%. If G0 makes a G2 cell unavailable, the substitute is the same $a$ at the nearest eligible $r$ toward 0.70 not already allocated, and a changed cell count requires recalibration before G2 (\code{simulation/v4_plan.py:substitute_cells}). This follows the amendment and the code, which exclude only cells already allocated; the earlier decision-specification wording, which excluded values of $r$ already in G2, is superseded by them. Before G2, a second binding record fixes the G0 status digest, the G2 cells and the criteria. It also applies a waste guard. The 54 G0 fits of B on the observed channel with $0<a\le0.30$ give a Jeffreys $\mathrm{Beta}(0.5,0.5)$ posterior for the technical-failure rate. From it, 200,000 beta-binomial draws with seed 250925006 estimate the probability of meeting the G2 and G2c allowances jointly. If that probability is below 0.80, execution stops before G2 (\code{simulation/v4_gate.py:g0_gap_binding} and \code{allocation_preflight}).

\subsubsection*{15. Seed namespaces}
% ODD36
Each generator seed is the first eight bytes, big-endian, of $\mathrm{SHA\text{-}256}$ of the string \code{namespace|block|phase|role|index}. The phases are \code{match}, \code{validate}, \code{train}, \code{check}, \code{reference}, \code{nuisance}, \code{bootstrap} and \code{stencil}. The roles are \code{state}, \code{null-state}, \code{network}, \code{mask}, \code{error} and \code{heap}. The block encodes repetition and sub-block (each 0--9,999) and the attempt (0 or 1) without aliases. The namespaces are:
\begin{itemize}
\item G0: \code{ic-detection-boundary-v4-calibration};
\item G1c, G2 and G2c: \code{ic-detection-boundary-v4-gate-rNNNN-calibration};
\item \code{ic-detection-boundary-v4-{validation,nuisance,bootstrap,stencil}-calibration};
\item confirmatory tiers: \code{ic-detection-boundary-v4-confirmatory} and \code{-{validation,nuisance,bootstrap,stencil}-confirmatory}, which the gate runner cannot derive.
\end{itemize}
Stencils reuse the base-phase streams to preserve common random numbers. In the tiers, one seed plan per value of $r$ is shared by all strengths and units at that $r$. Code: \code{simulation/detection_v4.py:SeedPlan.seed}, \code{simulation/v4_curve.py:ConfirmatorySeedPlan} and \code{simulation/v4_gate.py:task_seed_plan}.

\subsubsection*{16. Inventory, tier order, compute window and status categories}
% ODD37
The inventory has 46 comparison grids $\times$ 9 values of $r$ $\times$ 9 strengths, giving 3,726 cells in 414 comparison/$r$ units (\code{simulation/v4_plan.py:inventory}). The tiers run in fixed order, with a barrier between tiers:
\begin{enumerate}
\item B, observed;
\item A, C, D and E;
\item the three controls $\times$ five nulls, and mobility in three channels $\times$ five nulls;
\item $c=0.25/0.50$, $\delta=0.10/0.20$ and the exploratory combination;
\item $b=0$ and $\omega\in\{0.25,0.5,0.75,1\}$;
\item 24 steps per wave.
\end{enumerate}
Controls wait for their observed unit. At most eight deterministic shards run under the same code hash, on one machine whose environment matches the freeze record, and AUC, band width or invalidity never reorder or accelerate work. The window closes at 2026-10-18 23:59 America/Sao\_Paulo (2026-10-19T02:59:00Z). At the close, incomplete work is interrupted and preserved. Each unit is recorded as \emph{not started}, \emph{partial at operational close}, \emph{complete invalid} or \emph{complete valid} (recorded labels: \textit{não iniciado}, \textit{parcial por encerramento operacional}, \textit{completo inválido}, \textit{completo válido}). The linked amendment of 28 September authorizes release either when all 414 units have final curve records and no curve worker is running, or at the original deadline, without partial early release. % ODD37 RC47
Only complete units are released, together with the full inventory. % ODD37
Infrastructure errors suspend execution and allow resumption; they never become gaps. Code: \code{simulation/v4_curve.py:schedule}, \code{deadline}, \code{unit_status} and \code{released_curves}.

\subsection{Implementation details}

% ODD38
The model is implemented in Python 3.12 with NumPy, in the locked study package \code{ic_study}. The freeze record binds the interpreter, the NumPy version, the lock file, the platform, single-threaded BLAS and a cross-host numerical probe. The gate and curve runners refuse to execute unless a preflight reproduces every bound hash (\code{simulation/v4_gate.py:preflight}). Gate and confirmatory outputs are written to a sealed store outside any Git checkout, which also captures worker output (\code{simulation/v4_store.py:SealedStore}); before authorized release the runners print only gate summaries (statuses, counts, the per-fit matching diagnostics of G0, criteria values and pass/fail) and progress counts, and confirmatory curves are printed only by the release step (\code{run_gate.py} and \code{run_curve.py} in \code{experiments/exploratory/detection-boundary-v4/}).

% ODD40 RC48 RC49 RC50
The post-freeze release wrapper and G1c serialization successor are documented with their revisions, hashes and loader checks in \hyperref[supp:execution-records]{Supplementary Material: execution and reporting records} and \hyperref[supp:protocol-lineage]{the supplementary protocol genealogy and exposure inventory}. They preserve the frozen scientific code and gate criterion. The frozen \code{closed_utc} records the conventional deadline; \code{release-record.json} records actual time, trigger and output hash.
